Los sistemas de salud dependen de una logística que el paciente rara vez ve. Para que un examen se procese, un tratamiento se administre o un pabellón funcione, alguien tiene que llevar a tiempo una muestra al laboratorio, un medicamento refrigerado a la farmacia o un equipo de vuelta a su centro de reposición. Cuando ese transporte falla, el costo no es solo logístico: una muestra que llega tarde puede invalidar un examen y obligar a repetir la toma. En un caso de traslado de muestras biológicas entre centros de salud, cerca del 40\,\% de las muestras llegaba fuera de plazo con la planificación manual, y una planificación optimizada logró entregarlas todas a tiempo \cite{benini_detti_zabalo_2021}.

Este tipo de logística combina dificultades que en la literatura suelen estudiarse por separado. Cada envío debe recogerse en un lugar y entregarse en otro dentro de horarios acotados, lo que corresponde a los problemas de recogida y entrega \cite{parragh_doerner_hartl_2008}. La organización, además, decide con anticipación cuántos vehículos contratar y desde dónde operarlos, una decisión de localización que condiciona todos los recorridos posteriores \cite{prodhon_prins_2014}. Y cuando la flota propia no alcanza, el servicio puede encargarse a un tercero a un costo mayor, una alternativa que obliga a decidir qué conviene atender internamente \cite{dabia_lai_vigo_2019}. Planificar estas decisiones por partes, primero la flota y luego las rutas, es lo habitual en la práctica, pero puede dejar ahorros importantes sin capturar: en problemas de localización y ruteo, resolver ambas partes por separado ha llevado a soluciones bastante más caras que resolverlas en conjunto \cite{prodhon_prins_2014}.

El proyecto aborda este desafío para una red de transporte sanitario programado en la Región Metropolitana, que cada semana debe mover insumos, fármacos refrigerados, muestras clínicas y equipos entre más de cien establecimientos. La red debe atender cada solicitud dentro de su ventana y con recursos propios siempre que sea posible, al menor costo: equilibrando el gasto fijo de vehículos y bases con el de operarlos y el de recurrir a terceros. El informe propone una metodología que toma estas decisiones de forma coordinada y la compara con la forma en que un planificador las tomaría con reglas simples.

El Capítulo~\ref{cap:descripcion} describe el problema y los datos, el~\ref{cap:metodologia} la metodología, el~\ref{cap:casobase} el caso base y el~\ref{cap:trabajo} la implementación y los resultados en las instancias BIG y SMALL, incluida su variabilidad. El Capítulo~\ref{cap:sensibilidad} analiza la sensibilidad, el~\ref{cap:implementacion} evalúa alternativas de infraestructura y propone cómo implementar la herramienta con un piloto, y el~\ref{cap:conclusiones} concluye.
